% Two-Column Preprint -- a two-column article with a full-width title block.
% The page design (grid, fonts, accent colour, headings, header) is in
% preprint.sty; this file holds the content and the title metadata.
% Build with: latexmk -pdf main.tex
\documentclass[10pt, twoside]{article}
\usepackage{preprint}

\usepackage{amsmath}
\usepackage{pgfplots}
\pgfplotsset{compat=1.18}
\usepgfplotslibrary{groupplots}
\usetikzlibrary{positioning, arrows.meta}
\usepackage{siunitx}
\sisetup{mode=text, detect-weight, group-minimum-digits=4}

% Numeric references, compressed into ranges in the text.
\usepackage[backend=biber, style=numeric-comp, sorting=none, maxbibnames=4,
  giveninits=true, doi=false, url=false, isbn=false]{biblatex}
\addbibresource{refs.bib}
\renewcommand*{\bibfont}{\small}

\usepackage[hidelinks]{hyperref}
\usepackage[capitalise, noabbrev]{cleveref}

% Shared look for every plot: greys for baselines, the accent for our model.
\pgfplotsset{every axis/.append style={
  font=\sffamily\scriptsize, axis line style={black!45}, tick style={black!45},
  grid=major, grid style={black!8},
  legend style={draw=none, fill=none, font=\sffamily\scriptsize},
  every axis plot/.append style={line width=.9pt}}}

\DeclareMathOperator*{\argmax}{arg\,max}

% ---- Title metadata --------------------------------------------------------
\title{Citation-Aware Dense Retrieval\\for Scientific Literature Search}
\shorttitle{Citation-Aware Dense Retrieval}
\preprintkind{Preprint \enspace\textbullet\enspace Information Retrieval}
\preprintdate{October 2026}
\authors{Mira Okafor\affil{1,*}, Tomás Lindgren\affil{1}, Hana Sato\affil{2}
  and Daniel Ruiz-Brandt\affil{1,2}}
\affiliations{\affil{1}Department of Computer Science, University of Utrecht,
  The Netherlands \quad \affil{2}Institute for Language Technology,
  Kyoto University, Japan}
\correspondence{\affil{*}Correspondence: \texttt{m.okafor@example.org}.
  Code and indices will be released on publication.}
\keywords{dense retrieval, scientific literature, citation graphs,
  contrastive learning, hard negatives}

\begin{abstract}
Researchers rarely search for papers with the vocabulary of the papers they
need. Lexical rankers such as BM25 reward term overlap and miss work phrased in
the conventions of a neighbouring field, while general-purpose dense retrievers
are trained on web questions and transfer poorly to technical prose. We present
\textsc{Cadre}, a dual-encoder retriever that treats the citation graph as
supervision: co-cited papers become positives, and papers that share
vocabulary but are never cited together become hard negatives. On three
literature-search benchmarks \textsc{Cadre} improves nDCG@10 by 4.1 points
over the strongest dense baseline and recovers 91\% of the relevant papers in
the top 100, at the latency of a single nearest-neighbour lookup.
\end{abstract}

\begin{document}
\maketitle

\section{Introduction}
\label{sec:intro}

A literature search is a recall problem with a precision budget. A reader
drafting a related-work section wants every paper that bears on the claim, but
will read only the first page of results. Classical rankers built on the
probabilistic relevance framework~\cite{robertson2009bm25,manning2008ir}
remain strong on this task because scientific queries are full of rare,
discriminative terms. They fail when the same idea is named differently:
\emph{contrastive pre-training} in one community is \emph{metric learning} in
another, and a query in one dialect ranks the other far below the fold.

Dense retrievers map queries and documents into a shared vector space and
score them by inner product~\cite{halvorsen2021dual}. Their quality depends
almost entirely on the training pairs. Web-scale question--passage pairs teach
a model what a good answer to a factoid question looks like, not what a
relevant prior study looks like. Scientific text, however, comes with its own
relevance judgements: authors cite what they consider related. We exploit this
signal directly, building on earlier work that used citation contexts as weak
supervision for document embeddings~\cite{duval2020citation}.

Our contributions are threefold. First, we define positives and negatives
from the citation graph in a way that avoids the false negatives that plague
mined hard examples~\cite{oyelaran2022hardneg}. Second, we combine the
graph-derived loss with distillation from a cross-encoder~\cite{brennan2023distill}.
Third, we evaluate on three benchmarks drawn from twelve
fields~\cite{castellano2023scibench} and report where the model still fails
(\cref{sec:results}).

% Full-width floats are placed one page ahead of where they appear.
\begin{table*}[t]
  \centering
  \caption{Retrieval quality on the three SciRank splits. nDCG@10 and
  Recall@100 in percent; latency in milliseconds per query. The best score in
  each column is set in bold.}
  \label{tab:main}
  \small
  \begin{tabular}{@{}l S[table-format=2.1] S[table-format=2.1]
      S[table-format=2.1] S[table-format=2.1] S[table-format=2.1]
      S[table-format=2.1] S[table-format=3.0]@{}}
    \toprule
    & \multicolumn{2}{c}{Bio} & \multicolumn{2}{c}{Phys}
    & \multicolumn{2}{c}{CS} & \\
    \cmidrule(lr){2-3}\cmidrule(lr){4-5}\cmidrule(lr){6-7}
    Model & {nDCG@10} & {R@100} & {nDCG@10} & {R@100}
          & {nDCG@10} & {R@100} & {ms} \\
    \midrule
    BM25                        & 38.2 & 71.4 & 41.0 & 74.9 & 36.5 & 69.8 & 18 \\
    Dense, web pairs            & 35.7 & 76.3 & 37.9 & 77.2 & 39.1 & 80.4 & 9 \\
    Dense, mined negatives      & 40.8 & 81.0 & 42.6 & 82.5 & 43.3 & 84.7 & 9 \\
    Late interaction            & 43.5 & 84.2 & 44.9 & 85.1 & \bfseries 46.8 & 87.9 & 112 \\
    \midrule
    \textsc{Cadre} (ours)       & \bfseries 47.9 & \bfseries 90.6 & \bfseries 48.7
      & \bfseries 91.2 & 46.2 & \bfseries 91.4 & 9 \\
    \bottomrule
  \end{tabular}
\end{table*}

\begin{figure*}[t]
  \centering
  \begin{tikzpicture}
    \begin{groupplot}[group style={group size=3 by 1, horizontal sep=12mm},
        width=0.345\textwidth, height=44mm, xmode=log, log basis x=10,
        xmin=1, xmax=1000, ymin=0, ymax=100,
        xlabel={cut-off $k$}, ylabel={Recall@$k$ (\%)}]
      \nextgroupplot[title={Bio}]
        \addplot[black!35, dashed] coordinates {(1,9) (5,26) (10,37) (20,48) (50,62) (100,71) (200,79) (500,86) (1000,90)};
        \addplot[black!60] coordinates {(1,11) (5,31) (10,43) (20,55) (50,71) (100,81) (200,87) (500,92) (1000,95)};
        \addplot[accent] coordinates {(1,14) (5,38) (10,52) (20,65) (50,81) (100,91) (200,94) (500,97) (1000,98)};
      \nextgroupplot[title={Phys}, ylabel={}]
        \addplot[black!35, dashed] coordinates {(1,10) (5,28) (10,40) (20,51) (50,65) (100,75) (200,82) (500,88) (1000,91)};
        \addplot[black!60] coordinates {(1,12) (5,33) (10,45) (20,57) (50,73) (100,83) (200,88) (500,93) (1000,95)};
        \addplot[accent] coordinates {(1,15) (5,39) (10,53) (20,66) (50,82) (100,91) (200,95) (500,97) (1000,98)};
      \nextgroupplot[title={CS}, ylabel={},
          legend to name=recall-legend, legend columns=3,
          legend style={/tikz/every even column/.append style={column sep=4mm}}]
        \addplot[black!35, dashed] coordinates {(1,8) (5,24) (10,35) (20,46) (50,60) (100,70) (200,78) (500,85) (1000,89)};
        \addplot[black!60] coordinates {(1,12) (5,34) (10,47) (20,59) (50,76) (100,85) (200,90) (500,94) (1000,96)};
        \addplot[accent] coordinates {(1,14) (5,37) (10,51) (20,64) (50,81) (100,91) (200,95) (500,97) (1000,98)};
        \legend{BM25, Dense with mined negatives, \textsc{Cadre}}
    \end{groupplot}
  \end{tikzpicture}

  \vspace{2pt}
  \pgfplotslegendfromname{recall-legend}
  \caption{Recall as a function of the cut-off $k$ on a logarithmic axis.
  The gap between \textsc{Cadre} and the lexical baseline is widest between
  $k = 20$ and $k = 100$, the range a reader or re-ranker actually inspects.}
  \label{fig:recall}
\end{figure*}

\section{Method}
\label{sec:method}

\subsection{Dual-encoder scoring}

Let $q$ be a query and $d$ a document, represented by its title and abstract.
Two encoders $f_\theta$ and $g_\phi$ with shared weights map them to unit
vectors in $\mathbb{R}^{k}$, and the relevance score is their inner product,
\begin{equation}
  s(q, d) = \frac{f_\theta(q)^{\top} g_\phi(d)}{\tau},
  \label{eq:score}
\end{equation}
where $\tau > 0$ is a learned temperature. Because $g_\phi(d)$ does not depend
on the query, every document is encoded once and stored in an approximate
nearest-neighbour index~\cite{varga2024ann}; retrieval then costs one encoder
pass and one index lookup.

\subsection{Citation-derived training pairs}

For a paper $p$ with reference list $R(p)$, we treat two papers as related when
they are cited together by at least $m$ later papers. The positive set of a
query paper $p$ is
\begin{equation}
  P(p) = \bigl\{\, d \;:\; \lvert C(p) \cap C(d) \rvert \ge m \,\bigr\},
  \label{eq:positives}
\end{equation}
where $C(\cdot)$ is the set of papers citing its argument. Hard negatives are
drawn from the top BM25 results for the title of $p$, excluding any paper
within two hops of $p$ in the citation graph. This exclusion is what keeps
lexically similar but genuinely related work out of the negative set.

\subsection{Training objective}

Given a batch of $B$ queries, each with one positive $d^{+}_i$ and a set
$N_i$ of hard negatives, we minimise the contrastive loss
\begin{equation}
  \mathcal{L}_{\text{c}} = -\frac{1}{B}\sum_{i=1}^{B}
  \log \frac{e^{s(q_i, d^{+}_i)}}
  {e^{s(q_i, d^{+}_i)} + \sum_{d \in N_i \cup D^{-}_i} e^{s(q_i, d)}},
  \label{eq:loss}
\end{equation}
where $D^{-}_i$ contains the positives of the other queries in the batch. We
add a distillation term $\mathcal{L}_{\text{kd}}$ that matches the score
distribution of a cross-encoder teacher on the same candidates, and train on
$\mathcal{L} = \mathcal{L}_{\text{c}} + \lambda\,\mathcal{L}_{\text{kd}}$ with
$\lambda = 0.5$. \Cref{fig:arch} summarises the pipeline.

\begin{figure}[t]
  \centering
  \begin{tikzpicture}[font=\sffamily\scriptsize, node distance=5mm and 7mm,
      box/.style={draw=black!40, rounded corners=2pt, minimum height=7mm,
        minimum width=19mm, align=center, fill=white},
      key/.style={box, draw=accent, fill=accentpanel},
      arr/.style={-{Stealth[length=4pt]}, black!55, line width=.6pt}]
    \node[box] (q) {Query\\title or prompt};
    \node[box, below=of q] (d) {Document\\title + abstract};
    \node[key, right=of q] (fq) {Encoder $f_\theta$};
    \node[key, right=of d] (gd) {Encoder $g_\phi$};
    \node[box, right=11mm of $(fq)!.5!(gd)$] (s) {Score\\$s(q,d)$};
    \node[box, below=of gd, minimum width=30mm] (g) {Citation graph\\positives and negatives};
    \draw[arr] (q) -- (fq);
    \draw[arr] (d) -- (gd);
    \draw[arr] (fq) -- (s);
    \draw[arr] (gd) -- (s);
    \draw[arr, accent] (g.east) -| (s.south);
    \draw[black!40, dashed] (fq) -- node[right, black!55] {shared} (gd);
  \end{tikzpicture}
  \caption{Training pipeline. Both encoders share weights; the citation graph
  supplies the positive and negative documents for each query.}
  \label{fig:arch}
\end{figure}

\section{Experimental setup}
\label{sec:setup}

\paragraph{Corpora}
We pre-train on 9.6 million papers with their reference lists and evaluate on
the three SciRank splits~\cite{castellano2023scibench}: \emph{Bio} (life
sciences), \emph{Phys} (physics and astronomy) and \emph{CS} (computer
science). Each query is a paragraph from a related-work section, and the
relevant documents are the papers it cites.

\paragraph{Baselines}
We compare against BM25, a general dense retriever trained on web
questions~\cite{halvorsen2021dual}, the same retriever with mined hard
negatives~\cite{oyelaran2022hardneg}, and a late-interaction model that stores
one vector per token~\cite{sorensen2022lateinter}.

\paragraph{Metrics}
We report nDCG@10 for the first page and Recall@100 for the candidate set a
reader or a re-ranker would see. Latency is the median over \num{10000}
queries on one GPU, including the index lookup.


\section{Results}
\label{sec:results}

\Cref{tab:main} reports the main comparison. \textsc{Cadre} is the best model
on five of the six quality columns and matches the latency of the other
single-vector retrievers. Its largest gains are on Bio and Phys, where
vocabulary differs most between subfields; on CS, the late-interaction model
is 0.6 points ahead in nDCG@10 but answers twelve times more slowly and needs
an index roughly forty times larger.

\Cref{fig:recall} shows recall across cut-offs. At $k = 1$ all models are
close, because the top result is usually the paper whose title repeats the
query. The gap opens as $k$ grows: the papers that \textsc{Cadre} recovers and
BM25 misses are the ones phrased in another field's terms.

\paragraph{Ablation}
\Cref{tab:ablation} removes one component at a time on the Bio split. The
two-hop exclusion in \cref{eq:positives} matters most: without it, a quarter
of the mined negatives are papers the query actually cites, and the
contrastive loss in \cref{eq:loss} pushes them away.

\begin{table}[t]
  \centering
  \caption{Ablation on the Bio split. Each row removes one component from the
  full model.}
  \label{tab:ablation}
  \small
  \begin{tabular}{@{}l S[table-format=2.1] r@{}}
    \toprule
    Variant & {nDCG@10} & $\Delta$ \\
    \midrule
    Full model                     & 47.9 & -- \\
    No two-hop exclusion           & 43.1 & $-4.8$ \\
    No distillation ($\lambda = 0$) & 45.6 & $-2.3$ \\
    In-batch negatives only        & 44.0 & $-3.9$ \\
    Co-citation threshold $m = 1$  & 46.4 & $-1.5$ \\
    \bottomrule
  \end{tabular}
\end{table}

\paragraph{Failure cases}
Most remaining errors involve very recent papers. A paper published after the
citation snapshot has no co-citations, so its embedding is shaped only by its
text; on queries whose relevant set is dominated by such papers, nDCG@10 drops
to the level of the mined-negative baseline. Refreshing the graph monthly
closes most of this gap at modest cost, since only the encoder for new
documents needs to run.

\section{Conclusion}
\label{sec:conclusion}

Citations are the most abundant relevance judgements in science, and they come
for free. Using them to choose positives, and the citation graph to keep false
negatives out of the training batch, gives a single-vector retriever that
outperforms both lexical and general dense baselines on literature search
without a latency cost. The obvious next step is to weight citations by the
sentence they appear in, so that a citation of a method and a citation of a
contradicting result no longer count the same.

\paragraph{Acknowledgements}
We thank the maintainers of the open citation snapshots used in this work and
the anonymous reviewers of an earlier draft for their comments.

\printbibliography

\end{document}
